\chapter{Worked Examples and Exam Map}
\label{ch:exammap}
This chapter is a \emph{map}. It does not develop anything new; it shows how the
derivations of Parts~II--IV get used on exam-style questions, and where to go back to
when a question does not yield. Rather than isolated arithmetic, it works one circuit
(\cref{sec:onecircuit}) and one antenna (\cref{sec:oneantenna}) through everything the
exam asks about them---both taken unchanged from the chapters that derived them, so
every number here can be traced back. \Cref{ch:practice} then reasons through real pool items
one at a time, and \cref{ch:crossproblems} poses problems that need several chapters at
once.

\section{How the Topics Map to Exam-Style Thinking}
\begin{tabularx}{\textwidth}{@{}L{0.24\textwidth}Y@{}}
\toprule
Topic & What to recognize \\
\midrule
Units & Convert prefixes, distinguish \(f\) from \(\omega\), and check dimensions. \\
Complex impedance & Use rectangular form for series addition and polar form for magnitude/phase. \\
Resonance & Ask series or parallel, then ask whether the question concerns current, impedance, voltage, or bandwidth. \\
Poles and the s-plane & Translate ringing, selectivity, and stability into pole distance from the \(\jj\omega\) axis. \\
Filters & Identify passband, stopband, cutoff, slope, ripple, and bandwidth as pole/zero placement. \\
Feedback and stability & Reduce to loop gain: predictable gain, gain--bandwidth trade, and phase margin. \\
Decibels & Use \(10\log\) for power and \(20\log\) for voltage or current ratios under equal impedance. \\
Transmission lines & Use \(\Gamma\), SWR, return loss, velocity factor, and wavelength. \\
Active circuits & Separate gain, feedback, bandwidth, distortion, and stability. \\
\bottomrule
\end{tabularx}

\section{One Circuit, Seven Questions}
\label{sec:onecircuit}
Isolated examples teach the arithmetic and hide the point. So take a single
circuit---\cref{ch:rlc}'s running series RLC, with
\[
L=\SI{10}{\micro\henry},
\qquad
C=\SI{220}{\pico\farad},
\qquad
R=\SI{5}{\ohm},
\]
and ask it everything an exam might. Every answer below is a formula derived earlier,
cited so it can be checked, and no new component values are introduced.

\paragraph{1. Reactance of one part.}
At \(f=\SI{14.2}{\mega\hertz}\), that capacitor alone has
\(X_C=1/\omega C=\SI{51}{\ohm}\) (\cref{ch:ac}), so \(Z_C=-\jj\SI{51}{\ohm}\). The
sign and the \(\omega=2\pi f\) conversion are where these questions are lost, not the
division.

\paragraph{2. Resonant frequency.}
\[
f_0=\frac{1}{2\pi\sqrt{LC}}
=\frac{1}{2\pi\sqrt{(10^{-5})(2.2\times10^{-10})}}
\approx\SI{3.393}{\mega\hertz}
\qquad(\cref{ch:rlc}).
\]
\(R\) does not appear: any resistance in the question is a distractor for this part.
At \(f_0\) both reactances are \SI{213}{\ohm} and cancel exactly---which is the check
worth doing, since \(X_L=X_C\) \emph{is} the definition of resonance.

\paragraph{3. Quality factor.}
\(Q_s=\omega_0L/R=213/5\approx42.6\) (\cref{ch:rlc}). Note it is
\emph{this} topology's formula; \cref{ch:rlcpar} inverts it, and
\cref{sec:seriesparallel} shows why the two are one quantity.

\paragraph{4. Bandwidth---and the band edges.}
\(BW=f_0/Q=\SI{79.6}{\kilo\hertz}\), and \cref{sec:halfpower} showed this is an
\emph{equality}, not an estimate. The same section locates the edges: they straddle
\(f_0\) geometrically, at \SI{3.354}{\mega\hertz} and \SI{3.433}{\mega\hertz}, whose
\emph{product} is \(f_0^2\). A question asking for the edges rather than the width is
testing that distinction.

\paragraph{5. Voltage across a reactive element.}
Driven with \SI{10}{\volt} RMS at resonance, \(Z=R\), so \(I=10/5=\SI{2}{\ampere}\)
and
\(V_L=I\,\omega_0L=\SI{426}{\volt}\)---exactly \(Q\) times the source
(\cref{ch:rlc}). This is why tuned-circuit capacitors need voltage ratings the
transmitter's output would never suggest.

\paragraph{6. Pole location, and ringing.}
\(\alpha=R/2L=\SI{2.5e5}{\per\second}\) and \(\omega_d\approx\omega_0\), so the poles
sit at \(-2.5\times10^5\pm\jj\,2.13\times10^7\) (\cref{ch:rlc}). The envelope decays
with \(\tau=1/\alpha=\SI{4}{\micro\second}\)---and \cref{sec:ringdown}'s
\(\tau=1/\pi BW\) gives \SI{4}{\micro\second} too, from the bandwidth alone. Two
routes, one number.

\paragraph{7. What changes if \(R\) doubles?}
\(f_0\) does not move at all. \(Q\) halves, \(BW\) doubles, the poles move directly
away from the \(\jj\omega\) axis, and the ringing shortens by half
(\cref{sec:polegeometry}). Questions of the form ``which of these changes when\dots''
are asking whether you know which parameter each quantity depends on.

\begin{controlsbox}
\textbf{Why one circuit answers seven questions.} Nothing above was a separate fact to
memorize. Reactance, resonance, \(Q\), bandwidth, magnification, pole location and
ring-down are seven readings of one second-order transfer function, and the exam
samples them at random. That is the whole thesis of this book applied to a test: learn
the object, not the seven questions. If a question resists, identify which \emph{view}
of the pole pair it is asking for---frequency, time, or complex plane---and
\cref{ch:fourviews} is the chapter that connects them.
\end{controlsbox}

\section{One Antenna, Three Questions}
\label{sec:oneantenna}
The same treatment for \cref{ch:lines}'s running antenna, a resonant dipole presenting
\(R_{\text{rad}}=\SI{35}{\ohm}\) to \SI{50}{\ohm} coax.

\paragraph{Reflection coefficient.}
\(\Gamma=(35-50)/(35+50)=-0.176\) (\cref{sec:reflectionsection}). The negative sign
says the reflected wave is inverted, because the load is \emph{below} \(Z_0\).

\paragraph{SWR and return loss.}
\(\mathrm{SWR}=(1+0.176)/(1-0.176)=1.43\), and return loss is
\(-20\log_{10}|\Gamma|=\SI{15.1}{\decibel}\). Both are the same number read two ways,
and \cref{sec:sparams} adds a third name for it: \(|S_{11}|\).

\paragraph{How much power comes back.}
\(|\Gamma|^2=0.031\), so \SI{3.1}{\percent} is reflected and \SI{96.9}{\percent}
delivered. Note what this is \emph{not}: it says nothing about how much is
\emph{radiated}, which is \cref{sec:antennaeff}'s \(\eta\) and cannot be measured from
the shack.

\begin{exambox}
Decibel questions are the other place a single conversion is tested repeatedly, and one
worked line covers them. A filter that attenuates voltage to \(0.25\) of its input, at
equal source and load impedance, is
\[
20\log_{10}(0.25)=-\SI{12.0}{\decibel},
\qquad\text{and}\qquad
10\log_{10}(0.25^2)=10\log_{10}(0.0625)=-\SI{12.0}{\decibel}.
\]
The two agree \emph{because} the impedances are equal (\cref{sec:decibels}); when they
are not, the \(20\log\) form is simply wrong, and that qualification is what such a
question is usually probing.
\end{exambox}

\section{Where Each Topic Is Developed}
For the reasoning behind any row of the table, see: units and components
\cref{ch:foundations}; series/parallel combining \cref{ch:resistive}; impedance,
admittance, and complex power \cref{ch:ac}; first-order circuits \cref{ch:rc,ch:rl};
resonance and \(Q\) \cref{ch:rlc,ch:rlcpar}; poles, zeros, and Bode
\cref{ch:splane}; feedback and stability margins \cref{ch:feedback}; filters and
matching \cref{ch:filters}; transmission lines and the Smith chart \cref{ch:lines};
active circuits and oscillators \cref{ch:active}; and measurement
\cref{ch:measurement}.
